\begin{center}  \Large{\textbf{ABSTRACT}}  \end{center}
\begin{center}  \Large{\textbf{DIPLOMA THESIS TITLE}}  \end{center}
\vspace{3pt}
\begin{center}  \large{\textbf{Nikolaos Gerontas: \hspace{50pt} Konstantinos Moustakas:}}  \end{center}
\vspace{10pt}
\noindent

Hand-driven interaction with three-dimensional content is most often
delivered through specialised hardware --- depth cameras, head-mounted
displays, instrumented gloves, or hand-held controllers --- each of which
raises the cost and effort of entry to the very interactions it enables.
This thesis asks whether a hand-driven, interactive 3D application can
instead be built from nothing more than a webcam and a commodity
\acrshort{cpu}, with no headset, no depth sensor, and no instrumented gloves.

The obstacle is one of representation. Google MediaPipe recovers 21
landmarks per hand in real time on commodity hardware, but those landmarks
are a flat list of points whose scale varies with camera distance, carry
only wrist-relative depth, and encode neither a joint hierarchy nor
anatomically valid ranges of motion. The \acrshort{mano} parametric hand
model supplies the missing structure --- a fixed kinematic tree of joint
rotations, independent of hand size and viewpoint --- and bridging the two
is the central problem of the work.

The bridge is an analytical inverse-kinematics solver that recovers the
pose in closed form: a single \acrshort{svd}-based point-set alignment for
the global orientation, followed by one swing-only traversal of the
kinematic tree. The landmarks are first smoothed with a 1\euro{} filter,
and the solved rotations, a detected gesture, and a placement anchor are
serialised as a \acrshort{json} datagram and streamed over \acrshort{udp}
to a Unity application. There, a shared interaction layer drives four scenes
that put the same pose to markedly different uses: a companion-and-grappling
bear, a vertex-by-vertex face deformer, an editable city, and a painting
tool whose drawings leave the application as point clouds for surface
reconstruction.

The pipeline runs entirely on the \acrshort{cpu}, with no \acrshort{gpu}
involvement at any stage. A latency evaluation over \num{20234} frames
measured a mean software latency of \SI{30.68}{\ms}, with a 99th percentile
of \SI{40}{\ms} --- comfortably below the \SI{60}{\ms} threshold above which
lag becomes perceptible to users. The system thus embodies a deliberate
trade of tracking accuracy for accessibility, delivering fine-grained,
two-handed, three-dimensional interaction at interactive latency from the
camera already built into a laptop.\\\\

\noindent
\vspace{0.3cm}
\\ \textit{keywords} - \textbf{
    monocular,
    MediaPipe,
    kinematics,
    MANO,
    Unity 3D
}
